\begin{thebibliography}{143}
\addcontentsline{toc}{chapter}{参考文献}
\raggedright
\bibitem[1]{fcplt-ref-1} M. Akveld and D. Salamon, Loops of Lagrangian submanifolds and pseudo-holomorphic discs, Geom. Funct. Anal. 11 (2001), 609–650.

\bibitem[2]{fcplt-ref-2} V. I. Arnol’d, On a characteristic class entering into conditions of quantization, Funct. Anal. Appl. 1 (1967), 1–13.

\bibitem[3]{fcplt-ref-3} —–, Some remarks on symplectic monodromy of Milnor fibrations, in The Floer Memorial Volume (H. Hofer, C. Taubes, A. Weinstein, and E. Zehnder, eds.), Progr. Math. 133, Birkhäuser, 1995, pp. 99–104.

\bibitem[4]{fcplt-ref-4} V. Arnol’d, V. Goryunov, O. Lyashko, and V. Vasiliev, Singularity theory. I, Encyclopaedia Math. Sci. 6, Springer-Verlag, 1998.

\bibitem[5]{fcplt-ref-5} V. I. Arnol’d, S. M. Gusein-Zade, and A. N. Varchenko, Singularities of differentiable maps, Birkhäuser, 1988.

\bibitem[6]{fcplt-ref-6} P. S.Aspinwall, Some navigation rules for D-brane monodromy, J. Math. Phys. 42 (2001), 5534–5552.

\bibitem[7]{fcplt-ref-7} M. Atiyah, Topological quantum field theories, Inst. Hautes Études Sci. Publ. Math. 68 (1989), 175–186.

\bibitem[8]{fcplt-ref-8} —–, On analytic surfaces with double points, Proc. Royal Soc. Ser. A 247 (1958), 237–244.

\bibitem[9]{fcplt-ref-9} M. Atiyah and I. Singer, The index of elliptic operators. V, Ann. of Math. 93 (1971), 139–149.

\bibitem[10]{fcplt-ref-10} M. Audin and J. Lafontaine (eds.), Holomorphic curves in symplectic geometry, Progr. Math. 117, Birkhäuser, 1994.

\bibitem[11]{fcplt-ref-11} D. Auroux, Monodromy invariants in symplectic topology, Preprint math.SG/0304113, 2003.

\bibitem[12]{fcplt-ref-12} —–, Symplectic 4-manifolds as branched coverings of CP2 , Invent. Math. 139 (2000), 551–602.

\bibitem[13]{fcplt-ref-13} —–, Symplectic maps to projective spaces and symplectic invariants, Turkish J. Math. 25 (2001), 1–42.

\bibitem[14]{fcplt-ref-14} D. Auroux, L. Katzarkov, and D. Orlov, Mirror symmetry for weighted projective planes and their noncommutative deformations, Preprint math.AG/0404281, 2004.

\bibitem[15]{fcplt-ref-15} —–, Mirror symmetry for Del Pezzo surfaces: Vanishing cycles and coherent sheaves, Invent. Math. 166 (2006), 537–582.

\bibitem[16]{fcplt-ref-16} D. Auroux, V. Munoz, and F. Presas, Lagrangian submanifolds and Lefschetz pencils, J. Symplectic Geom. 3 (2005), 171–219.

\bibitem[17]{fcplt-ref-17} D. Auroux and I. Smith, Lefschetz pencils, branched covers and symplectic invariants, Preprint math.SG/0401021, 2004.

\bibitem[18]{fcplt-ref-18} P. Balmer and M. Schlichting, Idempotent completion of triangulated categories, J. Algebra 236 (2001), 819–834.

\bibitem[19]{fcplt-ref-19} A. Beilinson, Coherent sheaves on Pn and problems in linear algebra, Functional Anal. Appl. 12 (1978), 214–216.

\bibitem[20]{fcplt-ref-20} A. Beilinson, J. Bernstein, and P. Deligne, Faisceaux pervers, Analyse et topologie sur les espaces singuliers (I), Astérisque 100, Soc. Math. France, 1983, pp. 7–171.

\bibitem[21]{fcplt-ref-21} J. Bernstein and V. Lunts, Equivariant sheaves and functors, Lecture Notes in Math. 1578, Springer-Verlag, 1994.

\bibitem[22]{fcplt-ref-22} P. Biran and K. Cieliebak, Lagrangian embeddings into subcritical Stein manifolds, Israel J. Math. 127 (2002), 221–244.

\bibitem[23]{fcplt-ref-23} J. Birman and H. Hilden, On the mapping class groups of closed surfaces as covering spaces, in Advances in the theory of Riemann surfaces (L. Ahlfors, L. Bers, H. Farkas, R. Gunning, I. Kra and H. Rauch, eds.), Ann. of Math Stud. 66 (1971), Princeton University Press, pp. 81–115.

\bibitem[24]{fcplt-ref-24} M. Bökstedt and A. Neeman, Homotopy limits in triangulated categories, Compositio Math. 86 (1993), 209–234.

\bibitem[25]{fcplt-ref-25} A. I. Bondal, Helices, representations of quivers, and Koszul algebras, [117], pp. 75–95.

\bibitem[26]{fcplt-ref-26} —–, Representations of associative algebras and coherent sheaves, Math. USSR Izv. 34 (1990), 23–42.

\bibitem[27]{fcplt-ref-27} A. I. Bondal and M. M. Kapranov, Representable functors, Serre functors, and mutations, Math. USSR Izv. 35 (1990), 519–541.

\bibitem[28]{fcplt-ref-28} —–, Enhanced triangulated categories, Math. USSR Sb. 70 (1991), 93–107.

\bibitem[29]{fcplt-ref-29} B. Booss-Bavnbeck and K. Wojciechowski, Elliptic boundary value problems for Dirac operators, Birkhäuser, 1993.

\bibitem[30]{fcplt-ref-30} Yu. Chekanov, Differential algebras of Legendrian links, Invent. Math. 150 (2002), 441–483.

\bibitem[31]{fcplt-ref-31} C.-H. Cho, Holomorphic discs, spin structures, and the Floer cohomology of the Clifford torus, Internat. Math. Res. Notices 35 (2004), 1803–1843.

\bibitem[32]{fcplt-ref-32} C.-H. Cho andY.-G. Oh, Floer cohomology and disc instantons of Lagrangian torus fibers in toric Fano manifolds, Asian J. Math. 10 (2006), 773–814.

\bibitem[33]{fcplt-ref-33} O. Cornea and F. Lalonde, Cluster homology, Electron. Res. Announc. Amer. Math. Soc. 12 (2006), 1–12.

\bibitem[34]{fcplt-ref-34} V. de Silva, The index of a family of pseudo-holomorphic disks, Preprint, 2000.

\bibitem[35]{fcplt-ref-35} V. de Silva, R. Gautschi, J. Robbin, and D. Salamon, Combinatorial Floer homology, Manuscript, 2001.

\bibitem[36]{fcplt-ref-36} P. Deligne et al. (eds.), Quantum fields and strings: a course for mathematicians, Amer. Math. Soc., 1999.

\bibitem[37]{fcplt-ref-37} S. K. Donaldson, Floer homology groups in Yang-Mills theory, Cambridge Tracts in Math. 147, Cambridge University Press, 2002.

\bibitem[38]{fcplt-ref-38} —–, Symplectic submanifolds and almost-complex geometry, J. Differential Geom. 44 (1996), 666–705.

\bibitem[39]{fcplt-ref-39} —–, Lefschetz pencils on symplectic manifolds, J. Differential Geom. 53 (1999), 205–236.

\bibitem[40]{fcplt-ref-40} —–, Polynomials, vanishing cycles, and Floer homology, in Mathematics: frontiers and perspectives, Amer. Math. Soc., 2000, pp. 55–64.

\bibitem[41]{fcplt-ref-41} Yu. Drozd, Tame and wild matrix problems, in Thirteen papers in algebra and number theory, Amer. Math. Soc. Transl. 128, Amer. Math. Soc., 1986, pp. 31–55.

\bibitem[42]{fcplt-ref-42} Ya. Eliashberg and M. Gromov, Convex symplectic manifolds, in Several complex variables, Proc. Symposia Pure Math. 52, Amer. Math. Soc., 1991, pp. 135–162.

\bibitem[43]{fcplt-ref-43} J. Etnyre, L. Ng, and J. Sabloff, Invariants of Legendrian knots and coherent orientations, J. Symplectic Geom. 1 (2002), 321–367.

\bibitem[44]{fcplt-ref-44} A. Floer, Morse theory for Lagrangian intersections, J. Differential Geom. 28 (1988), 513–547.

\bibitem[45]{fcplt-ref-45} —–, A relative Morse index for the symplectic action, Comm. Pure Appl. Math. 41 (1988), 393–407.

\bibitem[46]{fcplt-ref-46} —–, The unregularized gradient flow of the symplectic action, Comm. Pure Appl. Math. 41 (1988), 775–813.

\bibitem[47]{fcplt-ref-47} —–, Witten’s complex and infinite dimensional Morse theory, J. Differential Geom. 30 (1989), 207–221.

\bibitem[48]{fcplt-ref-48} A. Floer and H. Hofer, Coherent orientations for periodic orbit problems in symplectic geometry, Math. Z. 212 (1993), 13–38.

\bibitem[49]{fcplt-ref-49} A. Floer, H. Hofer, and D. Salamon, Transversality in elliptic Morse theory for the symplectic action, Duke Math. J. 80 (1995), 251–292.

\bibitem[50]{fcplt-ref-50} U. Frauenfelder, Gromov convergence of holomorphic discs, Diplomarbeit, ETH Zürich, 2001.

\bibitem[51]{fcplt-ref-51} K. Fukaya, Floer homology for 3-manifolds with boundary (abstract), in Topology, Geometry, and Field theory, World Scientific, 1992, pp. 1–21.

\bibitem[52]{fcplt-ref-52} —–, Morse homotopy, \ensuremath{A_\infty} -categories, and Floer homologies, in Proceedings of GARC workshop on Geometry and Topology (H. J. Kim, ed.), Seoul National University, 1993, pp. 1–102.

\bibitem[53]{fcplt-ref-53} —–, Floer homology for three-manifolds with boundary I, Preprint, 1997.

\bibitem[54]{fcplt-ref-54} —–, Morse homotopy and its quantization, in Geometric topology (Athens, GA, 1993) (W. H. Kazez, ed.), American Math. Society, 1997, pp. 409–440.

\bibitem[55]{fcplt-ref-55} —–, Floer homology and mirror symmetry. II, in Minimal surfaces, geometric analysis and symplectic geometry (Baltimore, MD, 1999), Math. Soc. Japan, 2002, pp. 31–127.

\bibitem[56]{fcplt-ref-56} —–, Deformation theory, homological algebra, and mirror symmetry, in Geometry and physics of branes (Como, 2001), IOP, 2003, pp. 121–209.

\bibitem[57]{fcplt-ref-57} —–, Mirror symmetry of abelian varieties and multi-theta functions, J. Algebraic Geom. 11 (2002), 393–512.

\bibitem[58]{fcplt-ref-58} K. Fukaya and Y.-G. Oh, Zero-loop open strings in the cotangent bundle and Morse homotopy, Asian J. Math. 1 (1998), 96–180.

\bibitem[59]{fcplt-ref-59} K. Fukaya, P. Seidel, and I. Smith, Exact Lagrangian submanifolds in simply-connected cotangent bundles, Invent. Math., to appear.

\bibitem[60]{fcplt-ref-60} K. Fukaya,Y.-G. Oh, H. Ohta, and K. Ono, Lagrangian intersection Floer theory – anomaly and obstruction, Preprint, first version, 2000; second expanded version, 2006.

\bibitem[61]{fcplt-ref-61} S. Gelfand and Yu. Manin, Methods of homological algebra, Springer-Verlag, 1996.

\bibitem[62]{fcplt-ref-62} E. Getzler and J. D. S. Jones, \ensuremath{A_\infty} -algebras and the cyclic bar complex, Illinois J. Math. 34 (1990), 256–283.

\bibitem[63]{fcplt-ref-63} R. Gompf, Handlebody construction of Stein surfaces, Ann. of Math. 148 (1998), 619–693.

\bibitem[64]{fcplt-ref-64} —–, Symplectic structures from Lefschetz pencils in high dimensions, in Proceedings of the Casson Fest, Geom. Topol. Publishers, 2004, pp. 267–290.

\bibitem[65]{fcplt-ref-65} A. Gorodentsev, Exceptional objects and mutations in derived categories, in [117], pp. 57–74.

\bibitem[66]{fcplt-ref-66} A. Gorodentsev and S. Kuleshov, Helix theory, Moscow Math. J. 4 (2004), 377–440.

\bibitem[67]{fcplt-ref-67} A. Gorodentsev and A. Rudakov, Exceptional vector bundles on projective spaces, Duke Math. J. 54 (1987), 115–130.

\bibitem[68]{fcplt-ref-68} P. Griffiths and J. Morgan, Rational homotopy theory and differential forms, Progr. Math. 16, Birkhäuser, 1981.

\bibitem[69]{fcplt-ref-69} V. Gugenheim, On a perturbation theory for the homology of the loop-space, J. Pure Appl. Algebra 25 (1982), 197–205.

\bibitem[70]{fcplt-ref-70} V. Gugenheim, L. Lambe, and J. Stasheff, Algebraic aspects of Chen’s twisting cochain, Illinois J. Math. 34 (1990), 485–502.

\bibitem[71]{fcplt-ref-71} D. Happel, Hochschild cohomology of finite-dimensional algebras, in Séminaire d’algèbre Paul Dubreil et Marie-Paul Malliavin, 39ème Année (Paris, 1987/1988), Lecture Notes in Math. 1404, Springer-Verlag, 1989, pp. 108–126.

\bibitem[72]{fcplt-ref-72} J. Harris and I. Morrison, Moduli of curves, Grad. Texts in Math. 187, Springer-Verlag, 1998.

\bibitem[73]{fcplt-ref-73} F. Harvey and H. B. Lawson, Finite volume flows and Morse theory, Ann. of Math. 153 (2001), 1–25.

\bibitem[74]{fcplt-ref-74} R. Hind, Lagrangian isotopies in Stein manifolds, Preprint math.SG/0311093, 2003.

\bibitem[75]{fcplt-ref-75} —–, Lagrangian spheres in \(S^2\times S^2\) , Geom. Funct. Anal. 14 (2004), 303–318.

\bibitem[76]{fcplt-ref-76} H. Hofer, Pseudoholomorphic curves in symplectizations with applications to the Weinstein conjecture in dimension three, Invent. Math. 114 (1993), 515–563.

\bibitem[77]{fcplt-ref-77} Y.-Z. Huang, Two-dimensional conformal geometry and vertex operator algebras, Progr. Math. 148, Birkhäuser, 1997.

\bibitem[78]{fcplt-ref-78} M. Hutchings, Reidemeister torsion in generalized Morse theory, Forum Math. 14 (2002), 209–244.

\bibitem[79]{fcplt-ref-79} E.-N. Ionel and T. Parker, Gromov-Witten invariants of symplectic sums, Math. Res. Lett. 5 (1998), 563–576.

\bibitem[80]{fcplt-ref-80} T. Kadeishvili, The algebraic structure in the homology of an A.1/-algebra, Soobshch. Akad. Nauk Gruzin. SSR 108 (2) (1982), 249–252.

\bibitem[81]{fcplt-ref-81} —–, The category of differential coalgebras and the category of A.1/-algebras, Trudy Tbiliss. Mat. Inst. Razmadze Akad. Nauk Gruzin. SSR 77 (1985), 50–70.

\bibitem[82]{fcplt-ref-82} B. Keller, Introduction to A-infinity algebras and modules, Homology Homotopy Appl. 3 (2001), 1–35.

\bibitem[83]{fcplt-ref-83} —–, Bimodule complexes via strong homotopy actions, Algebr. Represent. Theory 3 (2000), 357–376.

\bibitem[84]{fcplt-ref-84} —–, A-infinity algebras in representation theory, in Representations of algebras, vol. 1, Beijing Normal University Press, 2002, pp. 74–86.

\bibitem[85]{fcplt-ref-85} M. Khovanov and P. Seidel, Quivers, Floer cohomology, and braid group actions, J. Amer. Math. Soc. 15 (2002), 203–271.

\bibitem[86]{fcplt-ref-86} R. Kirby and L. Taylor, Pin structures on low-dimensional manifolds, in Geometry of low-dimensional manifolds (S. Donaldson and C. B. Thomas, eds.), vol. 2, London Math. Soc. Lecture Note Ser. 151, Cambridge University Press, 1990, pp. 177–242.

\bibitem[87]{fcplt-ref-87} F. Knudsen and D. Mumford, The projectivity of the moduli space of stable curves I: preliminaries on “det” and “div”, Math. Scand. 39 (1976), 19–55.

\bibitem[88]{fcplt-ref-88} M. Kontsevich, Deformation quantization of Poisson manifolds I, Lett. Math. Phys. 66 (2003), 157–216.

\bibitem[89]{fcplt-ref-89} —–, Lectures at ENS Paris, Spring 1998, set of notes taken by J. Bellaiche, J.-F. Dat, I. Marin, G. Racinet and H. Randriambololona.

\bibitem[90]{fcplt-ref-90} —–, Homological algebra of mirror symmetry, in Proceedings of the International Congress of Mathematicians (Zürich, 1994), Birkhäuser, 1995, pp. 120–139.

\bibitem[91]{fcplt-ref-91} D. Kwon andY.-G. Oh, Structure of the image of (pseudo)-holomorphic discs with totally real boundary conditions, Comm. Anal. Geom. 8 (2000), 31–82.

\bibitem[92]{fcplt-ref-92} K. Lamotke, The topology of complex projective varieties after S. Lefschetz, Topology 20 (1981), 15–51.

\bibitem[93]{fcplt-ref-93} H. B. Lawson and M.-L. Michelsohn, Spin geometry, Princeton University Press, 1989.

\bibitem[94]{fcplt-ref-94} Y.-J. Lee, Reidemeister torsion in symplectic Floer theory and counting pseudoholomorphic tori, I, J. Symplectic Geom. 3 (2005), 221–311.

\bibitem[95]{fcplt-ref-95} K. Lefevre, Sur les \ensuremath{A_\infty} -catégories, Ph.D. thesis, Université Paris 7, 2002.

\bibitem[96]{fcplt-ref-96} J.-L. Loday, Cyclic homology, Springer-Verlag, 2nd edition, 1997.

\bibitem[97]{fcplt-ref-97} V. Lyubashenko, Category of \ensuremath{A_\infty} -categories, Homology Homotopy Appl. 5 (2003), 1–48.

\bibitem[98]{fcplt-ref-98} M. Markl, Transferring \ensuremath{A_\infty} (strongly homotopy associative) structures, Rend. Circ. Mat. Palermo .2/ Suppl. No. 79 (2006), 139–151.

\bibitem[99]{fcplt-ref-99} —–, Homotopy algebras are homotopy algebras, Forum Math. 16 (2004), 129–160.

\bibitem[100]{fcplt-ref-100} D. McDuff, Symplectic manifolds with contact type boundaries, Invent. Math. 103 (1991), 651–671.

\bibitem[101]{fcplt-ref-101} D. Nadler, Microlocal branes are constructible sheaves, Preprint math.SG/0612399, 2007.

\bibitem[102]{fcplt-ref-102} D. Nadler and E. Zaslow, Constructible Sheaves and the Fukaya Category, Preprint math.SG/0604379, 2007.

\bibitem[103]{fcplt-ref-103} Y.-G. Oh, Floer cohomology of Lagrangian intersections and pseudo-holomorphic discs I, Comm. Pure Appl. Math. 46 (1993), 949–994.

\bibitem[104]{fcplt-ref-104} C. Okonek, M. Schneider, and H. Spindler, Vector bundles on complex projective spaces, Progr. Math. 3, Birkhäuser, 1980.

\bibitem[105]{fcplt-ref-105} P. Piazza, Determinant bundles, manifolds with boundary and surgery, Comm. Math. Phys. 178 (1996), 597–626.

\bibitem[106]{fcplt-ref-106} S. Piunikhin, D. Salamon, and M. Schwarz, Symplectic Floer-Donaldson theory and quantum cohomology, in Contact and symplectic geometry (C. B. Thomas, ed.), Cambridge University Press, 1996, pp. 171–200.

\bibitem[107]{fcplt-ref-107} A. Polishchuk, Massey and Fukaya products on elliptic curves, Adv. Theor. Math. Phys. 4 (2000), 1187–1207.

\bibitem[108]{fcplt-ref-108} A. Polishchuk and E. Zaslow, Categorical mirror symmetry: the elliptic curve, Adv. Theor. Math. Phys. 2 (1998), 443–470.

\bibitem[109]{fcplt-ref-109} L. Polterovich, Surgery of Lagrange submanifolds, Geom. Funct. Anal. 1 (1991), 198–210.

\bibitem[110]{fcplt-ref-110} S. Priddy, Koszul resolutions, Trans. Amer. Math. Soc. 152 (1970), 39–60.

\bibitem[111]{fcplt-ref-111} A. Prouté, Algèbres différentielles fortement homotopiquement associatives, Thèse d’État, Université Paris VII, 1985.

\bibitem[112]{fcplt-ref-112} J. Robbin and D. Salamon, The Maslov index for paths, Topology 32 (1993), 827–844.

\bibitem[113]{fcplt-ref-113} —–, The spectral flow and the Maslov index, Bull. London Math. Soc. 27 (1995), 1–33.

\bibitem[114]{fcplt-ref-114} —–, Asymptotic behaviour of holomorphic strips, Ann. Inst. H. Poincaré Anal. Non Linéaire 18 (2001), 573–612.

\bibitem[115]{fcplt-ref-115} Y. Ruan and G. Tian, A mathematical theory of Quantum Cohomology, J. Differential Geom. 42 (1995), 259–367.

\bibitem[116]{fcplt-ref-116} —–, Higher genus symplectic invariants and sigma models coupled with gravity, Invent. Math. 130 (1997), 455–516.

\bibitem[117]{fcplt-ref-117} A. Rudakov et al., Helices and vector bundles: Seminaire Rudakov, London Math. Soc. Lecture Note Ser. 148, Cambridge University Press, 1990.

\bibitem[118]{fcplt-ref-118} D. Salamon and E. Zehnder, Morse theory for periodic solutions of Hamiltonian systems and the Maslov index, Comm. Pure Appl. Math. 45 (1992), 1303–1360.

\bibitem[119]{fcplt-ref-119} M. Schwarz, Cohomology operations from S 1 -cobordisms in Floer homology, Ph.D. thesis, ETH Zürich, 1995.

\bibitem[120]{fcplt-ref-120} —–, A quantum cup-length estimate for symplectic fixed points, Invent. Math. 133 (1998), 353–397.

\bibitem[121]{fcplt-ref-121} S. Scott, Splitting the curvature of the determinant line bundle, Proc. Amer. Math. Soc. 128 (2000), 2763–2775.

\bibitem[122]{fcplt-ref-122} G. Segal, Lectures on topological field theory (Stanford notes), Manuscript, 1999.

\bibitem[123]{fcplt-ref-123} P. Seidel, \(\pi_1\) of symplectic automorphism groups and invertibles in quantum homology rings, Geom. Funct. Anal. 7 (1997), 1046–1095.

\bibitem[124]{fcplt-ref-124} —–, Lagrangian two-spheres can be symplectically knotted, J. Differential Geom. 52 (1999), 145–171.

\bibitem[125]{fcplt-ref-125} —–, Graded Lagrangian submanifolds, Bull. Soc. Math. France 128 (2000), 103–146.

\bibitem[126]{fcplt-ref-126} —–, More about vanishing cycles and mutation, in Symplectic Geometry and Mirror Symmetry (Proceedings of the 4th KIAS Annual International Conference) (K. Fukaya, Y.-G. Oh, K. Ono, and G. Tian, eds.), World Scientific, 2001, pp. 429–465.

\bibitem[127]{fcplt-ref-127} —–, Vanishing cycles and mutation, in Proceedings of the 3rd European Congress of Mathematics (Barcelona), Vol. 2, Birkhäuser, 2001, pp. 65–85.

\bibitem[128]{fcplt-ref-128} —–, Braids and symplectic four-manifolds with abelian fundamental group, Turkish J. Math. 26 (2001 Gököva Geometry and Topology Conference Proceedings Special Volume) (2002), 93–100.

\bibitem[129]{fcplt-ref-129} —–, Fukaya categories and deformations, in Proceedings of the International Congress of Mathematicians (Beijing, 2002), Vol. II, Higher Ed. Press, 2002, pp. 351–360.

\bibitem[130]{fcplt-ref-130} —–, Homological mirror symmetry for the quartic surface, Preprint math.SG/0310414, 2003.

\bibitem[131]{fcplt-ref-131} —–, A long exact sequence for symplectic Floer cohomology, Topology 42 (2003), 1003–1063.

\bibitem[132]{fcplt-ref-132} —–, Exact Lagrangian submanifolds in \(T^*S^n\) and the graded Kronecker quiver, in Different faces of geometry, Kluwer/Plenum, 2004, pp. 349–364.

\bibitem[133]{fcplt-ref-133} P. Seidel and I. Smith, The symplectic topology of Ramanujam’s surface, Comment. Math. Helv. 80 (2005), 859–881.

\bibitem[134]{fcplt-ref-134} P. Seidel and R. Thomas, Braid group actions on derived categories of coherent sheaves, Duke Math. J. 108 (2001), 37–108.

\bibitem[135]{fcplt-ref-135} V. De Silva, Products in the symplectic Floer homology of Lagrangian intersections, Ph.D. thesis, Oxford University, 1998.

\bibitem[136]{fcplt-ref-136} N. Spaltenstein, Resolutions of unbounded complexes, Compositio Math. 65 (1988), 121–154.

\bibitem[137]{fcplt-ref-137} J. Stasheff, Homotopy associativity of H -spaces I, Trans. Amer. Math. Soc. 108 (1963), 275–292.

\bibitem[138]{fcplt-ref-138} R. Switzer, Algebraic topology - homotopy and homology, Grundlehren Math. Wiss. 212, Springer-Verlag, 1975.

\bibitem[139]{fcplt-ref-139} R. Thomas, Stability conditions and the braid group, Comm. Anal. Geom. 14 (2006), 135–161.

\bibitem[140]{fcplt-ref-140} R. Thomas and S.-T. Yau, Special Lagrangians, stable bundles and mean curvature flow, Comm. Anal. Geom. 10 (2002), 1075–1113.

\bibitem[141]{fcplt-ref-141} J.-L. Verdier, Des catégories dérivées des catégories abéliennes, Astérisque 239, Soc. Math. France, 1996.

\bibitem[142]{fcplt-ref-142} B. Wajnryb, Artin groups and geometric monodromy, Invent. Math. 138 (1999), 563–571.

\bibitem[143]{fcplt-ref-143} C. A. Weibel, An introduction to homological algebra, Cambridge University Press, 1994.

\end{thebibliography}
